\documentclass[11pt,a4paper]{article}

\usepackage[utf8]{inputenc}
\usepackage[T1]{fontenc}
\usepackage{lmodern}
\usepackage[spanish]{babel}
\usepackage[margin=2.5cm]{geometry}
\usepackage{estilocrispropioLibro}

\pagestyle{fancy}
\fancyhf{}
\lhead{Física y Química --- 3\textsuperscript{o} ESO}
\rhead{\nouppercase{\rightmark}}
\cfoot{\thepage}

\title{\vspace{-1.5cm}\textbf{Situaciones de Aprendizaje}\\[4pt]\Large Física y Química --- 3\textsuperscript{o} ESO}
\author{}
\date{}

\begin{document}
\maketitle
\tableofcontents

%===================================================================
\newpage
\section{Unidad 1: Método Científico}
\subsection*{1. Título}
``¿Cuánta agua gastamos realmente? Investigando el consumo de agua en nuestro instituto con el método científico''

\subsection*{2. Justificación}
Esta situación de aprendizaje pretende que el alumnado aplique de forma práctica las etapas del método científico, el uso de magnitudes y unidades del Sistema Internacional, la notación científica y la representación gráfica de datos, trabajados en la Unidad 1, a un problema real y medible de su propio centro educativo. Se busca desarrollar la competencia científica mediante la indagación, el trabajo cooperativo y el uso crítico de fuentes digitales, conectando el aprendizaje con la sostenibilidad y el ahorro de recursos, valores especialmente relevantes en Castilla-La Mancha, una región con escasez estructural de agua.

\subsection*{3. Contextualización}
Castilla-La Mancha es una de las comunidades autónomas españolas con mayor riesgo de estrés hídrico y sequía recurrente (cuencas del Guadiana y del Júcar). Los centros educativos consumen agua a diario (baños, cafetería, limpieza, riego de zonas verdes) sin que el alumnado sea consciente de las cantidades reales ni de cómo se miden y expresan estas magnitudes. Planteamos un reto cercano: ``El AMPA y el equipo directivo del instituto quieren reducir la factura de agua y piden a las clases de 3º de ESO un informe científico sobre el consumo actual y propuestas de ahorro''. Este contexto conecta directamente con la vida diaria del alumnado y con un problema socioambiental real de su entorno geográfico.

\subsection*{4. Fundamentación curricular}
\textbf{Competencias específicas relacionadas}:
\begin{itemize}
    \item \textbf{CE1} (Comprender fenómenos fisicoquímicos): aplicar el concepto de magnitud, unidad y densidad al analizar el consumo de agua.
    \item \textbf{CE2} (Método científico e indagación): formular preguntas e hipótesis sobre el consumo de agua y comprobarlas con datos reales.
    \item \textbf{CE3} (Lenguaje científico): manejar unidades del SI (litros, metros cúbicos), notación científica y gráficas.
    \item \textbf{CE4} (Uso crítico de recursos digitales): buscar y contrastar información fiable sobre consumo de agua y recomendaciones de ahorro.
    \item \textbf{CE5} (Trabajo colaborativo): desarrollar el proyecto en equipos con roles definidos.
    \item \textbf{CE6} (Ciencia y sociedad): relacionar el problema con la sostenibilidad y las necesidades ambientales de Castilla-La Mancha.
\end{itemize}

\textbf{Criterios de evaluación trabajados}: CE1.1, CE1.2, CE1.3, CE2.1, CE2.2, CE3.1, CE3.2, CE4.1, CE4.2, CE5.1, CE5.2, CE6.2.

\textbf{Saberes básicos}:
\begin{itemize}
    \item \textbf{Bloque A (Destrezas científicas básicas)}: ``Metodologías de la investigación científica: identificación y formulación de cuestiones, elaboración de hipótesis y comprobación experimental de las mismas''; ``El lenguaje científico: unidades del Sistema Internacional y sus símbolos. Herramientas matemáticas básicas en diferentes escenarios científicos y de aprendizaje''; ``Estrategias de interpretación y producción de información científica utilizando diferentes formatos y diferentes medios''.
    \item \textbf{Bloque C (La energía)}: conexión transversal con ``Concienciación sobre la necesidad del ahorro energético y la conservación sostenible del medio ambiente'' (aplicado aquí al recurso agua como paralelismo de sostenibilidad).
\end{itemize}

\subsection*{5. Metodología}
Se emplea un enfoque de \textbf{aprendizaje basado en proyectos (ABP)} con estructura cooperativa. La clase se organiza en grupos heterogéneos de 3-4 estudiantes, cada uno con un rol rotativo y definido:
\begin{itemize}
    \item \textbf{Coordinador/a}: organiza el trabajo del grupo y controla los plazos.
    \item \textbf{Medidor/a de datos}: responsable de tomar y anotar las mediciones o estimaciones de consumo.
    \item \textbf{Investigador/a de fuentes}: busca y selecciona la información en internet.
    \item \textbf{Portavoz/redactor final}: coordina la redacción del informe final y la presentación oral.
\end{itemize}
Se combinan sesiones de trabajo en el aula (cálculo, notación científica, gráficas) con una sesión de ``trabajo de campo'' en el propio centro (observación de contadores o grifos, si es posible, o estimación con datos facilitados por conserjería/dirección) y con un tiempo de investigación digital guiada.

\subsection*{6. Recursos}
\begin{itemize}
    \item Datos reales o simulados del contador de agua del centro (facilitados por el equipo directivo o conserjería).
    \item Calculadora científica y hojas de cálculo (LibreOffice Calc o Google Sheets) para tablas y gráficas.
    \item Fuentes digitales fiables sugeridas: página del \textbf{Instituto Geológico y Minero de España (IGME)}, la \textbf{Confederación Hidrográfica del Guadiana} (organismo público de cuenca en Castilla-La Mancha), y el portal del \textbf{Canal de Isabel II} o \textbf{Aqualia} para datos de consumo doméstico medio. Buscador recomendado: Google con las palabras clave \textit{``consumo medio de agua litros persona día España sitio:.gob.es''} o \textit{``estrés hídrico Castilla-La Mancha datos oficiales''}, priorizando dominios \texttt{.gob.es}, \texttt{.es} institucionales o de organismos públicos frente a blogs sin autoría.
    \item Simulador \textbf{PhET ``Estimation'' o ``Density''} (opcional) para reforzar el manejo de magnitudes y unidades.
    \item Material de aula: cuaderno, ficha de recogida de datos, ordenador o tablet por grupo.
\end{itemize}
\textbf{Criterio de fiabilidad de fuentes}: solo se aceptan páginas de organismos oficiales (dominios .gob.es, .es institucionales, universidades o organizaciones científicas reconocidas), nunca blogs personales, foros o redes sociales sin autoría. \textbf{Formato mínimo de cita}: Autor o institución (Año). \textit{Título de la página}. Nombre del sitio web. URL completa. Fecha de consulta.

\subsection*{7. Tareas y actividades}
\begin{enumerate}
    \item \textbf{Sesión 1 (aula)}: presentación del reto por parte del profesorado, formación de grupos y reparto de roles. Lluvia de ideas sobre qué usos de agua existen en el instituto.
    \item \textbf{Sesión 2 (aula/campo)}: cada grupo recoge o recibe datos reales de consumo de agua del centro (litros/mes, litros/día, número de grifos, etc.) y los organiza en una tabla, expresando las cantidades grandes en notación científica y convirtiéndolas a las unidades SI adecuadas (litros a metros cúbicos) usando factores de conversión.
    \item \textbf{Sesión 3 (aula de informática)}: investigación digital guiada siguiendo el criterio de fiabilidad de fuentes indicado en el apartado 6, para comparar el consumo del instituto con la media española y detectar si hay estrés hídrico en la región.
    \item \textbf{Sesión 4 (aula)}: elaboración de una gráfica (consumo mensual o por franjas horarias) siguiendo las normas de construcción de gráficas trabajadas en la unidad (variable independiente/dependiente, escala, título, interpretación de proporcionalidad).
    \item \textbf{Sesión 5 (aula)}: redacción del informe final grupal y preparación de una breve presentación oral (5 minutos) con al menos una propuesta razonada de ahorro de agua para el centro.
\end{enumerate}

\textbf{Cuestiones que el alumnado debe responder y entregar (informe final individual o grupal)}:
\begin{enumerate}
    \item ¿Cuál es el consumo total de agua del instituto en el periodo estudiado, expresado en litros y en notación científica en metros cúbicos?
    \item Aplicando factores de conversión, ¿a cuántos litros por alumno y día equivale ese consumo?
    \item ¿La relación entre el número de días transcurridos y el volumen de agua consumido es directamente proporcional? Justifica tu respuesta a partir de la gráfica elaborada.
    \item ¿Qué fuentes digitales habéis utilizado para comparar el consumo del centro con la media española, y por qué las consideráis fiables? Cita al menos dos fuentes con el formato indicado.
    \item Formula una hipótesis sobre qué factor (número de alumnado, día de la semana, estación del año) más influye en el consumo de agua del centro, y describe cómo la comprobaríais mediante un pequeño experimento o recogida de datos adicional.
    \item Proponed y justificad, con cálculos numéricos de ahorro estimado (en litros y en €, si se dispone del precio del agua), al menos dos medidas concretas de reducción del consumo de agua en el instituto.
    \item ¿Qué relación tiene este problema con la situación hídrica de Castilla-La Mancha? Argumenta tu respuesta con los datos encontrados.
\end{enumerate}

\subsection*{8. Evaluación}
Se evaluará mediante una rúbrica que combine:
\begin{itemize}
    \item El \textbf{informe escrito grupal} (cálculos correctos, uso adecuado de unidades SI y notación científica, gráfica bien construida, fuentes citadas correctamente): criterios CE1.1, CE1.2, CE3.1, CE3.2, CE4.2.
    \item La \textbf{hipótesis formulada y el procedimiento de comprobación propuesto}: CE2.1, CE2.2.
    \item El \textbf{funcionamiento del equipo} (roles cumplidos, participación equitativa), observado durante las sesiones de aula: CE5.1, CE5.2.
    \item La \textbf{presentación oral final} (claridad, argumentación, conexión con el problema real de Castilla-La Mancha): CE1.3, CE6.2.
\end{itemize}
Se emplearán instrumentos variados: rúbrica de informe, lista de cotejo de trabajo en equipo y coevaluación entre grupos tras las presentaciones orales.

\subsection*{9. Observaciones}
Si no es posible obtener datos reales del contador del centro, el profesorado puede facilitar una tabla de datos simulados pero realistas (basados en medias oficiales) para no interrumpir la secuencia didáctica. Se recomienda coordinar esta situación de aprendizaje con el departamento de Biología y Geología o con el equipo de sostenibilidad del centro (si existe un proyecto ``Escuela Verde'' o similar) para dar continuidad al reto más allá de la unidad. Debe prestarse atención a la diversidad del alumnado, permitiendo adaptar el rol de ``medidor de datos'' a estudiantes con dificultades de cálculo mediante apoyo visual (tablas prellenadas) y ofreciendo plantillas de la gráfica para quienes lo necesiten.

%===================================================================
\newpage
\section{Unidad 2: Cinemática}
\subsection*{1. Título}
``Ruta segura al instituto: analizamos la cinemática de nuestros desplazamientos diarios''

\subsection*{2. Justificación}
Esta situación de aprendizaje busca que el alumnado aplique de forma significativa los conceptos de posición, desplazamiento, velocidad y aceleración trabajados en la Unidad 2 a una realidad que vive a diario: sus propios trayectos de casa al instituto (a pie, en bicicleta, en autobús escolar o en coche). Se pretende desarrollar la competencia científica mediante indagación con datos reales, el trabajo en equipo, el uso crítico de fuentes digitales (mapas, aplicaciones de rutas) y la reflexión sobre la seguridad vial, un tema de gran relevancia para adolescentes que empiezan a moverse de forma más autónoma.

\subsection*{3. Contextualización}
Muchos municipios de Castilla-La Mancha tienen centros educativos alejados de las zonas residenciales, con carreteras comarcales sin arcenes o aceras insuficientes, lo que hace que la seguridad vial en los trayectos escolares sea una preocupación real para familias y ayuntamientos. Planteamos el reto: ``El Ayuntamiento quiere mejorar la seguridad de las rutas escolares y ha pedido a los institutos un estudio del alumnado sobre las velocidades, tiempos y distancias de los trayectos habituales, para identificar los puntos más peligrosos y proponer mejoras''. Este contexto es cercano a la experiencia diaria del alumnado y conecta la física con un problema social real de su entorno.

\subsection*{4. Fundamentación curricular}
\textbf{Competencias específicas relacionadas}:
\begin{itemize}
    \item \textbf{CE1} (Comprender fenómenos fisicoquímicos): aplicar los conceptos de velocidad y aceleración a situaciones cotidianas de movimiento.
    \item \textbf{CE2} (Método científico e indagación): formular hipótesis sobre qué medio de transporte o ruta es más rápido/seguro y comprobarlo con datos.
    \item \textbf{CE3} (Lenguaje científico): manejar correctamente unidades SI de velocidad y aceleración, y representar/interpretar gráficas x-t y v-t.
    \item \textbf{CE4} (Uso crítico de recursos digitales): emplear aplicaciones de mapas y buscadores para obtener distancias y tiempos fiables.
    \item \textbf{CE5} (Trabajo colaborativo): desarrollar el estudio en equipos con roles definidos.
    \item \textbf{CE6} (Ciencia y sociedad): relacionar los resultados con la seguridad vial y la mejora de infraestructuras del entorno.
\end{itemize}

\textbf{Criterios de evaluación trabajados}: CE1.1, CE1.2, CE1.3, CE2.1, CE2.2, CE3.1, CE3.2, CE4.1, CE4.2, CE5.1, CE5.2, CE6.2.

\textbf{Saberes básicos}:
\begin{itemize}
    \item \textbf{Bloque D (La interacción)}: ``Predicción de movimientos sencillos a partir de los conceptos de la cinemática, formulando hipótesis comprobables sobre valores futuros de estas magnitudes, validándolas a través del cálculo numérico, la interpretación de gráficas o el trabajo experimental''; ``Aplicación de las leyes de Newton... y predecir los efectos de estas en situaciones cotidianas y de seguridad vial'' (conexión introductoria con la seguridad vial, que se ampliará en la Unidad 3).
    \item \textbf{Bloque A (Destrezas científicas básicas)}: ``El lenguaje científico: unidades del Sistema Internacional y sus símbolos''; ``Estrategias de interpretación y producción de información científica utilizando diferentes formatos y diferentes medios''.
\end{itemize}

\subsection*{5. Metodología}
Se utiliza un enfoque de \textbf{indagación guiada} combinado con \textbf{trabajo cooperativo}. La clase se organiza en grupos de 3-4 estudiantes con roles definidos:
\begin{itemize}
    \item \textbf{Cronometrador/a-medidor/a}: recoge o estima tiempos y distancias de los trayectos del grupo.
    \item \textbf{Calculador/a}: realiza las conversiones de unidades y los cálculos de velocidad y aceleración.
    \item \textbf{Investigador/a digital}: obtiene datos de distancia y tiempo mediante aplicaciones de mapas.
    \item \textbf{Portavoz/redactor final}: coordina la elaboración del informe y la presentación final.
\end{itemize}
Se combinan actividades de recogida de datos reales (autoinforme de cada estudiante sobre su trayecto), trabajo de cálculo y representación gráfica en el aula, e investigación digital guiada para contrastar los datos personales con información oficial de seguridad vial.

\subsection*{6. Recursos}
\begin{itemize}
    \item Aplicación de mapas (Google Maps o similar) para medir distancias y tiempos estimados de los trayectos de cada estudiante.
    \item Calculadora científica y hoja de cálculo (LibreOffice Calc o Google Sheets) para elaborar las gráficas x-t y v-t.
    \item Fuentes digitales fiables recomendadas: página de la \textbf{DGT (Dirección General de Tráfico)}, en especial su sección de estadísticas de accidentalidad y campañas de seguridad vial escolar, y el portal de la \textbf{Junta de Comunidades de Castilla-La Mancha} sobre movilidad escolar sostenible. Buscador recomendado: Google con las palabras clave \textit{``seguridad vial escolar estadísticas sitio:dgt.es''} o \textit{``camino escolar seguro Castilla-La Mancha''}, priorizando siempre dominios \texttt{.gob.es}, \texttt{.es} institucionales o ayuntamientos, frente a foros o blogs sin autoría verificable.
    \item Simulador \textbf{PhET ``Moving Man''} para comprobar de forma interactiva las gráficas de posición, velocidad y aceleración.
    \item Cronómetro (del móvil o de laboratorio) y cinta métrica, si se realiza la comprobación experimental en el patio.
\end{itemize}
\textbf{Criterio de fiabilidad de fuentes}: se aceptan únicamente páginas de organismos oficiales (DGT, ayuntamientos, Junta de Castilla-La Mancha, universidades), nunca foros, blogs personales o redes sociales sin autoría. \textbf{Formato mínimo de cita}: Autor o institución (Año). \textit{Título de la página}. Nombre del sitio web. URL completa. Fecha de consulta.

\subsection*{7. Tareas y actividades}
\begin{enumerate}
    \item \textbf{Sesión 1 (aula)}: presentación del reto, formación de grupos y reparto de roles. Cada estudiante recoge individualmente el tiempo aproximado y la distancia de su trayecto habitual al instituto (usando la app de mapas).
    \item \textbf{Sesión 2 (aula)}: puesta en común dentro del grupo; se calcula la velocidad media de cada trayecto individual, convirtiendo todas las unidades al SI, y se clasifica el tipo de movimiento (¿es razonable suponerlo URM?).
    \item \textbf{Sesión 3 (aula de informática)}: investigación digital guiada según el criterio de fiabilidad de fuentes del apartado 6, para comparar las velocidades medias obtenidas con los límites de velocidad legales según el medio de transporte (a pie, bici, coche) y con datos oficiales de accidentalidad en trayectos escolares.
    \item \textbf{Sesión 4 (patio, si es posible)}: comprobación experimental por parejas de la velocidad media al caminar y correr una distancia conocida, cronometrando el tiempo, siguiendo un protocolo de seguridad acordado previamente.
    \item \textbf{Sesión 5 (aula)}: elaboración de las gráficas x-t y v-t del grupo, redacción del informe final y preparación de una presentación oral (5 minutos) con una propuesta de mejora de seguridad vial para el trayecto más peligroso identificado.
\end{enumerate}

\textbf{Cuestiones que el alumnado debe responder y entregar (informe final grupal)}:
\begin{enumerate}
    \item Para cada miembro del grupo, calcula la velocidad media de su trayecto habitual en m/s, mostrando la conversión de unidades paso a paso.
    \item ¿Qué tipo de movimiento (URM o UARM) se ajusta mejor a cada trayecto? Justifica tu respuesta.
    \item Dibuja la gráfica posición-tiempo de al menos dos trayectos distintos del grupo y compáralos: ¿cuál tiene mayor pendiente y qué significa eso?
    \item A partir de las fuentes oficiales consultadas, ¿cuáles son los límites de velocidad y las recomendaciones de seguridad para peatones, ciclistas o vehículos en zonas escolares? Cita al menos dos fuentes con el formato indicado.
    \item Formula una hipótesis sobre qué trayecto o medio de transporte del grupo es más seguro (no solo más rápido) y explica cómo la comprobaríais con datos adicionales.
    \item Identifica el punto o tramo del trayecto que consideráis más peligroso y proponed una mejora concreta (ej. paso de peatones, reductor de velocidad, carril bici), justificándola con los conceptos de velocidad y aceleración estudiados.
    \item ¿Qué relación tiene este estudio con la seguridad vial en las carreteras y caminos escolares de vuestro municipio o de Castilla-La Mancha en general?
\end{enumerate}

\subsection*{8. Evaluación}
Se evaluará mediante una rúbrica que combine:
\begin{itemize}
    \item El \textbf{informe escrito grupal} (cálculos correctos, uso de unidades SI, gráficas bien construidas, fuentes citadas correctamente): criterios CE1.1, CE1.2, CE3.1, CE3.2, CE4.2.
    \item La \textbf{hipótesis y el diseño de comprobación experimental} de la sesión de patio: CE2.1, CE2.2.
    \item El \textbf{funcionamiento del equipo} (cumplimiento de roles, cooperación), observado durante todas las sesiones: CE5.1, CE5.2.
    \item La \textbf{presentación oral final} y la propuesta de mejora de seguridad vial: CE1.3, CE6.2.
\end{itemize}
Se emplearán instrumentos variados: rúbrica de informe, lista de cotejo de trabajo en equipo, y coevaluación entre grupos tras las presentaciones.

\subsection*{9. Observaciones}
Si algún alumno o alumna no puede realizar el trayecto a pie o en bicicleta (por ejemplo, va en autobús o coche desde una zona rural alejada), puede aportar igualmente sus datos, ya que enriquece la comparación entre medios de transporte. Debe supervisarse con especial cuidado la sesión práctica en el patio (protocolo de seguridad, zona despejada). Se recomienda coordinar esta situación de aprendizaje con el departamento de Tecnología o con la Policía Local/Municipal si el centro tiene programas de ``camino escolar seguro'', para dar mayor autenticidad y posible continuidad al reto.

%===================================================================
\newpage
\section{Unidad 3: Dinámica y Fuerzas}
\subsection*{1. Título}
``Diseñamos una atracción de feria segura: fuerzas, muelles y seguridad en los parques de Castilla-La Mancha''

\subsection*{2. Justificación}
Esta situación de aprendizaje pretende que el alumnado aplique de manera integrada los conceptos de fuerza, suma de fuerzas, Ley de Hooke y Leyes de Newton, trabajados en la Unidad 3, al diseño y análisis de seguridad de una atracción de feria sencilla (por ejemplo, un columpio con muelles o un sistema de frenado por resorte). Se busca desarrollar la competencia científica mediante el trabajo experimental, la indagación con datos reales, el trabajo en equipo y el uso crítico de fuentes digitales, conectando la física con un contexto lúdico y motivador que además tiene implicaciones reales de seguridad para el alumnado y sus familias.

\subsection*{3. Contextualización}
Cada año, numerosos municipios de Castilla-La Mancha celebran ferias y fiestas patronales con atracciones mecánicas (columpios, norias, camas elásticas) que dependen de principios físicos de fuerzas, muelles y equilibrio para funcionar de forma segura. Planteamos el reto: ``Una empresa de atracciones de feria que trabaja en las fiestas de nuestra comarca quiere contratar a un equipo de jóvenes 'ingenieros' para revisar, mediante cálculos físicos, si el sistema de muelles de una atracción cumple los requisitos mínimos de seguridad, y quiere un informe técnico''. Este contexto conecta la física con una experiencia festiva y cercana al alumnado, aportando a la vez un componente de responsabilidad y seguridad.

\subsection*{4. Fundamentación curricular}
\textbf{Competencias específicas relacionadas}:
\begin{itemize}
    \item \textbf{CE1} (Comprender fenómenos fisicoquímicos): aplicar los conceptos de fuerza, Ley de Hooke y Leyes de Newton al análisis de un sistema mecánico real.
    \item \textbf{CE2} (Método científico e indagación): formular hipótesis sobre la relación fuerza-elongación y comprobarla experimentalmente con un muelle.
    \item \textbf{CE3} (Lenguaje científico): manejar correctamente unidades SI de fuerza (N) y constante elástica (N/m), interpretar y producir gráficas fuerza-elongación, y respetar las normas de seguridad del laboratorio.
    \item \textbf{CE4} (Uso crítico de recursos digitales): buscar y contrastar información fiable sobre normativa de seguridad de atracciones y parques.
    \item \textbf{CE5} (Trabajo colaborativo): desarrollar el proyecto en equipos con roles definidos.
    \item \textbf{CE6} (Ciencia y sociedad): relacionar los resultados con la seguridad pública y las tradiciones festivas de Castilla-La Mancha.
\end{itemize}

\textbf{Criterios de evaluación trabajados}: CE1.1, CE1.2, CE1.3, CE2.1, CE2.2, CE3.1, CE3.2, CE3.3, CE4.1, CE4.2, CE5.1, CE5.2, CE6.2.

\textbf{Saberes básicos}:
\begin{itemize}
    \item \textbf{Bloque D (La interacción)}: ``Las fuerzas como agentes de cambio: relación de los efectos de las fuerzas, tanto en el estado de movimiento o de reposo de un cuerpo como produciendo deformaciones en los sistemas sobre los que actúan''; ``Aplicación de las leyes de Newton: observación de situaciones cotidianas o de laboratorio que permiten entender cómo se comportan los sistemas materiales ante la acción de las fuerzas y predecir los efectos de estas en situaciones cotidianas y de seguridad vial''; ``Fenómenos gravitatorios, eléctricos y magnéticos: experimentos sencillos que evidencian la relación con las fuerzas de la naturaleza''.
    \item \textbf{Bloque A (Destrezas científicas básicas)}: ``Trabajo experimental y proyectos de investigación: estrategias en la resolución de problemas y en el desarrollo de investigaciones mediante la indagación, la deducción, la búsqueda de evidencias y el razonamiento lógico-matemático''; ``Normas de uso de cada espacio, asegurando y protegiendo así la salud propia y comunitaria''.
\end{itemize}

\subsection*{5. Metodología}
Se emplea un enfoque de \textbf{aprendizaje basado en problemas (ABP)} con trabajo experimental de laboratorio. La clase se organiza en grupos de 3-4 estudiantes con roles definidos:
\begin{itemize}
    \item \textbf{Responsable de laboratorio}: gestiona el material y verifica el cumplimiento de las normas de seguridad.
    \item \textbf{Medidor/a}: toma las medidas de elongación del muelle con la regla.
    \item \textbf{Investigador/a digital}: busca información sobre normativa de seguridad de atracciones.
    \item \textbf{Portavoz/redactor final}: elabora el informe técnico final y la presentación.
\end{itemize}
Se combinan sesiones de trabajo experimental en el aula-laboratorio (montaje de un muelle con pesas conocidas, siguiendo la Ley de Hooke) con sesiones de cálculo, representación gráfica e investigación digital guiada sobre normativa real de seguridad de atracciones.

\subsection*{6. Recursos}
\begin{itemize}
    \item Material de laboratorio: muelle o resorte, soporte, juego de pesas calibradas, regla o cinta métrica, dinamómetro.
    \item Calculadora científica y hoja de cálculo (LibreOffice Calc o Google Sheets) para tablas y gráficas fuerza-elongación.
    \item Fuentes digitales fiables recomendadas: página del \textbf{Ministerio de Industria, Comercio y Turismo} (normativa de seguridad de atracciones de feria en España, Real Decreto de aparatos de elevación y transporte), y el portal de la \textbf{Junta de Comunidades de Castilla-La Mancha} sobre normativa de ferias y espectáculos públicos. Buscador recomendado: Google con las palabras clave \textit{``normativa seguridad atracciones feria España sitio:.gob.es''} o \textit{``inspección técnica atracciones feriales Castilla-La Mancha''}, priorizando siempre dominios \texttt{.gob.es}, \texttt{.es} institucionales o ayuntamientos frente a foros o blogs sin autoría.
    \item Simulador \textbf{PhET ``Hooke's Law'' (``Masses and Springs'')} para reforzar de forma interactiva la relación fuerza-elongación.
    \item Vídeos divulgativos sobre el funcionamiento físico de atracciones de feria (columpios, norias), citados correctamente.
\end{itemize}
\textbf{Criterio de fiabilidad de fuentes}: se aceptan únicamente páginas de organismos oficiales (ministerios, Junta de Castilla-La Mancha, ayuntamientos, universidades) o simuladores educativos reconocidos como PhET; nunca foros, blogs personales o redes sociales sin autoría. \textbf{Formato mínimo de cita}: Autor o institución (Año). \textit{Título de la página}. Nombre del sitio web. URL completa. Fecha de consulta.

\subsection*{7. Tareas y actividades}
\begin{enumerate}
    \item \textbf{Sesión 1 (aula)}: presentación del reto, formación de grupos y reparto de roles. Repaso de la Ley de Hooke y formulación por cada grupo de una hipótesis sobre la relación entre la fuerza aplicada y la elongación del muelle.
    \item \textbf{Sesión 2 (laboratorio)}: montaje del experimento: cada grupo cuelga pesas conocidas de un muelle, mide la elongación correspondiente y anota los datos en una tabla, respetando las normas de seguridad del laboratorio (manejo cuidadoso de pesas, no soltar el muelle bruscamente, etc.).
    \item \textbf{Sesión 3 (aula)}: elaboración de la gráfica fuerza-elongación con los datos obtenidos y cálculo de la constante elástica $k$ del muelle a partir de la pendiente de la recta.
    \item \textbf{Sesión 4 (aula de informática)}: investigación digital guiada según el criterio de fiabilidad de fuentes indicado en el apartado 6, sobre la normativa real de seguridad de atracciones de feria, para comparar los requisitos técnicos con los resultados de su propio experimento.
    \item \textbf{Sesión 5 (aula)}: redacción del informe técnico final y preparación de una presentación oral (5 minutos) en la que el grupo, como ``ingenieros consultores'', explique si el muelle estudiado sería adecuado para una atracción real y qué fuerza máxima podría soportar con seguridad.
\end{enumerate}

\textbf{Cuestiones que el alumnado debe responder y entregar (informe final grupal)}:
\begin{enumerate}
    \item Con los datos obtenidos en el laboratorio, dibuja la gráfica fuerza-elongación de vuestro muelle y calcula su constante elástica $k$, indicando las unidades.
    \item ¿Los datos obtenidos confirman o refutan la hipótesis inicial sobre la relación entre fuerza y elongación? Justifica tu respuesta.
    \item Si la atracción real necesitara soportar una fuerza máxima de 500 N sin superar una elongación de 0.8 m, ¿sería adecuado vuestro muelle? Justifica el cálculo.
    \item A partir de las fuentes oficiales consultadas, ¿qué requisitos mínimos de seguridad exige la normativa española o autonómica para las atracciones de feria? Cita al menos dos fuentes con el formato indicado.
    \item Explica, aplicando la Primera Ley de Newton, qué le ocurriría a un pasajero de la atracción si el sistema de frenado (muelle) fallara repentinamente.
    \item Si dos personas de distinto peso se sientan en el mismo columpio con muelles, ¿cómo afectaría esto a la fuerza que deben soportar los muelles? Aplica la Segunda Ley de Newton y el concepto de peso ($p = m \cdot g$) en tu explicación.
    \item Como conclusión del informe técnico, ¿recomendaríais o no el uso del muelle estudiado para la atracción? Argumentad la respuesta con los cálculos y la normativa consultada.
\end{enumerate}

\subsection*{8. Evaluación}
Se evaluará mediante una rúbrica que combine:
\begin{itemize}
    \item El \textbf{informe de laboratorio} (hipótesis, procedimiento, tabla y gráfica de datos, cálculo correcto de $k$, unidades SI): criterios CE1.1, CE1.2, CE2.1, CE2.2, CE3.1, CE3.2.
    \item El \textbf{cumplimiento de las normas de seguridad} durante la práctica experimental, observado directamente: CE3.3.
    \item El \textbf{funcionamiento del equipo} (roles cumplidos, cooperación) durante todas las sesiones: CE5.1, CE5.2.
    \item El \textbf{informe técnico final y la presentación oral}, incluyendo el uso correcto de fuentes citadas y la argumentación sobre seguridad: CE1.3, CE4.2, CE6.2.
\end{itemize}
Se emplearán instrumentos variados: rúbrica de informe de laboratorio, lista de cotejo de seguridad, lista de cotejo de trabajo en equipo y coevaluación entre grupos tras las presentaciones.

\subsection*{9. Observaciones}
Es imprescindible supervisar de cerca la sesión de laboratorio para garantizar un uso seguro del material (pesas, muelles) y adaptar las pesas utilizadas a las características del muelle disponible en el centro. Si no se dispone de material de laboratorio suficiente, se puede sustituir el muelle real por una goma elástica calibrada o por el simulador PhET, manteniendo la misma secuencia de hipótesis-experimentación-análisis. Se recomienda coordinar esta situación de aprendizaje con el departamento de Tecnología (para el diseño técnico) o con el AMPA/Ayuntamiento si el centro tiene relación con las fiestas patronales locales, para dar mayor autenticidad al reto.

%===================================================================
\newpage
\section{Unidad 4: Energía}
\subsection*{1. Título}
``¿Cuánto cuesta encender el instituto? Auditoría energética y propuesta de ahorro para nuestro centro''

\subsection*{2. Justificación}
Esta situación de aprendizaje pretende que el alumnado aplique los conceptos de energía, trabajo y potencia trabajados en la unidad a un problema real y cercano: el consumo energético de su propio centro educativo. Se busca desarrollar la competencia científica mediante el cálculo de magnitudes energéticas reales, fomentar la conciencia sobre el ahorro energético y la sostenibilidad, y practicar el trabajo cooperativo y el uso crítico de fuentes de información, contribuyendo así al desarrollo competencial integral del alumnado más allá del mero cálculo numérico.

\subsection*{3. Contextualización}
Castilla-La Mancha es una de las comunidades autónomas con mayor potencial e implantación de energía eólica y solar en España, pero también sufre veranos muy calurosos e inviernos fríos que disparan el consumo energético en los centros educativos (calefacción, refrigeración, iluminación, equipos informáticos). El alumnado convive a diario con interruptores de luz, radiadores, ordenadores y proyectores en las aulas sin ser consciente de su coste energético y económico real. Esta situación aprovecha ese entorno cotidiano --- el propio instituto --- como contexto de aprendizaje, conectando el consumo energético del centro con la factura eléctrica familiar que ya se analiza en la unidad.

\subsection*{4. Fundamentación curricular}
\textbf{Competencias específicas relacionadas}:
\begin{itemize}
    \item \textbf{CE1} --- Comprender y explicar fenómenos fisicoquímicos del entorno (el consumo energético de aparatos eléctricos) y resolver problemas aplicando las leyes de la energía y la potencia.
    \item \textbf{CE3} --- Manejar con soltura unidades de medida (J, W, kWh) y la interpretación de datos e información (facturas, fichas técnicas de aparatos).
    \item \textbf{CE4} --- Utilizar de forma crítica y eficiente plataformas digitales y recursos variados para consultar información fiable.
    \item \textbf{CE5} --- Utilizar estrategias de trabajo colaborativo para desarrollar un proyecto científico con impacto en la comunidad educativa.
    \item \textbf{CE6} --- Comprender la ciencia como herramienta para dar solución a necesidades ambientales y sociales (el ahorro energético y la sostenibilidad del centro).
\end{itemize}

\textbf{Criterios de evaluación concretos}: CE1.2, CE1.3, CE3.1, CE3.2, CE4.2, CE5.1, CE5.2, CE6.2.

\textbf{Saberes básicos}:
\begin{itemize}
    \item \textbf{Bloque C (La energía)}: ``Diseño y comprobación experimental de hipótesis relacionadas con el uso doméstico e industrial de la energía en sus distintas formas y las transformaciones entre ellas.''
    \item \textbf{Bloque C (La energía)}: ``Naturaleza eléctrica de la materia: electrización de los cuerpos, circuitos eléctricos y la obtención de energía eléctrica. Concienciación sobre la necesidad del ahorro energético y la conservación sostenible del medio ambiente.''
    \item \textbf{Bloque C (La energía)}: ``Elaboración fundamentada de hipótesis sobre el medio ambiente y la sostenibilidad a partir de las diferencias entre fuentes de energía renovables y no renovables. Valoración de la producción de energía eólica en Castilla-La Mancha.''
    \item \textbf{Bloque A (Las destrezas científicas básicas)}: ``El lenguaje científico: unidades del Sistema Internacional y sus símbolos. Herramientas matemáticas básicas en diferentes escenarios científicos y de aprendizaje.''
\end{itemize}

\subsection*{5. Metodología}
Se emplea un enfoque de aprendizaje basado en proyectos (ABP) con trabajo cooperativo estructurado. El alumnado se organiza en grupos heterogéneos de 3-4 personas, con roles rotativos y definidos:
\begin{itemize}
    \item \textbf{Coordinador/a}: organiza el reparto de tareas y los plazos del grupo.
    \item \textbf{Investigador/a de datos}: recopila las potencias de los aparatos y las tarifas eléctricas.
    \item \textbf{Responsable de cálculo}: realiza y verifica todos los cálculos de energía y coste.
    \item \textbf{Portavoz/diseñador/a}: elabora el informe final o cartel y lo presenta al resto de la clase.
\end{itemize}
Se combina el trabajo de campo (observación directa de aparatos en el aula), la búsqueda documental guiada y la puesta en común final en gran grupo, favoreciendo la autonomía, la interacción verbal y la integración del conocimiento matemático y científico.

\subsection*{6. Recursos}
\begin{itemize}
    \item Ficha de recogida de datos (plantilla en papel o digital) para anotar aparatos eléctricos del aula, su potencia (en la etiqueta o placa de características) y horas de uso diario estimadas.
    \item Ordenadores o tablets del aula/dispositivo personal para la búsqueda de información.
    \item \textbf{Fuentes de información concretas a utilizar:}
    \begin{itemize}
        \item Web de la \textbf{Comisión Nacional de los Mercados y la Competencia (CNMC)} o del \textbf{IDAE} (Instituto para la Diversificación y Ahorro de la Energía) para el precio medio del kWh en España.
        \item Web de \textbf{Red Eléctrica de España (REE)}, sección ``Demanda y generación en tiempo real'', para consultar el porcentaje de energía renovable en la red en el momento de la consulta.
        \item Buscador general (Google) con palabras clave sugeridas: \textit{``precio kWh España 2026 IDAE''}, \textit{``potencia contratada instituto consumo eléctrico''}, \textit{``energía eólica Castilla-La Mancha datos''}.
        \item Una factura eléctrica real (aportada voluntariamente por alguna familia, con los datos económicos ocultos) o la factura de ejemplo ya incluida en la unidad del libro.
    \end{itemize}
    \item \textbf{Criterio de fiabilidad de fuentes:} se dará prioridad a páginas institucionales (dominios .gob.es, .es de organismos oficiales) o de empresas eléctricas reconocidas, evitando blogs sin autoría o páginas sin fecha de actualización. Cada grupo deberá justificar por qué considera fiable la fuente elegida (autoría, organismo, fecha).
    \item \textbf{Formato mínimo de cita:} Nombre de la organización o autor (si consta). \textit{Título de la página}. Nombre del sitio web. URL completa. Fecha de consulta (día/mes/año).
    \item Calculadora científica o aplicación de calculadora.
\end{itemize}

\subsection*{7. Tareas y actividades}
\begin{enumerate}
    \item \textbf{Fase 1 (individual/grupo, 1 sesión):} Cada grupo elige un aula del instituto (puede ser la propia) y recoge datos in situ: número y tipo de aparatos eléctricos (luces, ordenador, proyector, radiador/aire acondicionado si procede), su potencia en vatios (leída en la etiqueta o buscada en internet a partir del modelo) y una estimación razonada de las horas de uso diario.
    \item \textbf{Fase 2 (grupo, 1 sesión):} Búsqueda documental guiada del precio del kWh actual en España y del porcentaje de energía renovable en la red eléctrica en el momento de la consulta, utilizando las fuentes indicadas en el apartado de Recursos, con cita correcta de cada fuente.
    \item \textbf{Fase 3 (grupo, 1-2 sesiones):} Cálculo de la energía consumida (en J y en kWh) por cada aparato al día, a la semana y en un trimestre, y del coste económico aproximado asociado, aplicando las fórmulas $E = P \cdot t$ y el precio del kWh obtenido.
    \item \textbf{Fase 4 (grupo, 1 sesión):} Elaboración de un informe o cartel final con los resultados y al menos dos propuestas concretas y viables de ahorro energético para el aula estudiada (por ejemplo, apagar proyectores en los descansos, aprovechar la luz natural, ajustar el termostato).
    \item \textbf{Fase 5 (gran grupo, 1 sesión):} Puesta en común: cada grupo presenta sus resultados y propuestas al resto de la clase (5 minutos por grupo), y se elabora una lista conjunta de recomendaciones para el centro.
\end{enumerate}

\textbf{Cuestiones que el alumnado deberá responder y entregar por escrito (individual o grupal) al finalizar:}
\begin{enumerate}
    \item ¿Cuál es la potencia total (en W y en kW) de los aparatos eléctricos del aula estudiada?
    \item ¿Cuánta energía consume el aula en un día lectivo completo, expresada en julios y en kWh?
    \item ¿Cuál es el coste económico aproximado de ese consumo diario, semanal y trimestral, según el precio del kWh consultado (indicando la fuente y fecha de consulta)?
    \item ¿Qué aparato es responsable de la mayor parte del consumo energético del aula? Justifica tu respuesta con los cálculos realizados.
    \item ¿Qué porcentaje de la energía eléctrica consumida en España (o en la red en el momento de vuestra consulta) procede de fuentes renovables? ¿Qué papel juega Castilla-La Mancha en la producción de energía eólica según lo que habéis encontrado?
    \item Proponed y justificad al menos dos medidas concretas de ahorro energético para el aula estudiada, estimando (de forma aproximada) cuánta energía y dinero se ahorrarían al aplicarlas.
    \item Reflexión grupal: ¿cómo os habéis repartido el trabajo dentro del grupo? ¿Qué dificultades habéis encontrado al buscar y contrastar la información de las distintas fuentes?
\end{enumerate}

\subsection*{8. Evaluación}
Se evaluará mediante una rúbrica que contemple: (a) la corrección de los cálculos de energía, potencia y coste (ligado a CE1.2 y CE3.2); (b) la calidad y fiabilidad de las fuentes utilizadas, con citación correcta (ligado a CE3.1 y CE4.2); (c) la viabilidad y fundamentación científica de las propuestas de ahorro energético (ligado a CE1.3 y CE6.2); (d) el funcionamiento cooperativo del grupo, valorado a través de la reflexión final escrita y la observación directa del profesorado durante las sesiones de trabajo (ligado a CE5.1 y CE5.2); y (e) la calidad de la presentación oral final ante el grupo clase. Se combinará la evaluación del producto final (informe/cartel) con la observación continua del proceso de trabajo en el aula.

\subsection*{9. Observaciones}
Es recomendable coordinar esta situación de aprendizaje con el equipo directivo o el/la responsable de mantenimiento del centro, que puede aportar datos reales de consumo o incluso visitar el aula si el centro dispone de contador inteligente. Se puede ampliar la actividad invitando a los grupos a presentar sus propuestas de ahorro más viables al equipo directivo, dando así una proyección real al trabajo del alumnado. Conviene adaptar el nivel de exigencia de la búsqueda de fuentes según la diversidad del grupo, ofreciendo enlaces ya preseleccionados al alumnado que lo necesite.

%===================================================================
\newpage
\section{Unidad 5: Sistemas Materiales}
\subsection*{1. Título}
``El laboratorio del suero casero: disoluciones que salvan vidas''

\subsection*{2. Justificación}
Esta situación de aprendizaje pretende que el alumnado comprenda y aplique de forma significativa los conceptos de concentración (porcentaje en masa, porcentaje en volumen y concentración en masa) trabajados en la unidad, a través de un contexto sanitario cercano y motivador: la preparación de sueros de rehidratación oral, un producto real que se usa en botiquines domésticos y centros de salud. Se busca desarrollar la competencia matemática aplicada a la química, el pensamiento crítico ante la información sanitaria y el trabajo cooperativo, conectando la ciencia escolar con un problema de salud pública reconocible por el alumnado.

\subsection*{3. Contextualización}
En los meses de verano, con las altas temperaturas propias de Castilla-La Mancha, son frecuentes los casos de deshidratación, especialmente en personas mayores y en niños y niñas pequeños, así como en la práctica deportiva. En los centros de salud y farmacias se recomienda el uso de sueros de rehidratación oral, cuya composición y concentración de sales y azúcares está cuidadosamente calculada. El alumnado ha podido ver (o usar) estos sobres de suero en casa o conoce a alguien que los ha necesitado. Esta situación aprovecha ese contexto realista para trabajar el cálculo de concentraciones con un producto de uso cotidiano y relevancia sanitaria.

\subsection*{4. Fundamentación curricular}
\textbf{Competencias específicas relacionadas}:
\begin{itemize}
    \item \textbf{CE1} --- Comprender y resolver problemas fisicoquímicos relacionados con la composición de disoluciones, aplicándolos a la mejora de la calidad de vida (salud).
    \item \textbf{CE2} --- Formular hipótesis y comprobar experimentalmente aspectos relacionados con la preparación de disoluciones.
    \item \textbf{CE3} --- Manejar con soltura las unidades de concentración y la interpretación de datos de etiquetas y prospectos.
    \item \textbf{CE4} --- Utilizar de forma crítica recursos digitales para consultar información sanitaria fiable.
    \item \textbf{CE5} --- Trabajar de forma colaborativa en la elaboración de un proyecto con utilidad para la comunidad (folleto informativo).
\end{itemize}

\textbf{Criterios de evaluación concretos}: CE1.1, CE1.2, CE2.2, CE3.1, CE3.2, CE4.2, CE5.1.

\textbf{Saberes básicos}:
\begin{itemize}
    \item \textbf{Bloque B (La materia)}: ``Experimentos relacionados con los sistemas materiales: conocimiento y descripción de sus propiedades, su composición y su clasificación. Técnicas de separación de mezclas.''
    \item \textbf{Bloque A (Las destrezas científicas básicas)}: ``Trabajo experimental y proyectos de investigación: estrategias en la resolución de problemas y en el desarrollo de investigaciones mediante la indagación, la deducción, la búsqueda de evidencias y el razonamiento lógico-matemático.''
    \item \textbf{Bloque A (Las destrezas científicas básicas)}: ``El lenguaje científico: unidades del Sistema Internacional y sus símbolos. Herramientas matemáticas básicas en diferentes escenarios científicos y de aprendizaje.''
    \item \textbf{Bloque A (Las destrezas científicas básicas)}: ``Diversos entornos y recursos de aprendizaje científico como el laboratorio o los entornos virtuales: materiales, sustancias y herramientas tecnológicas.''
\end{itemize}

\subsection*{5. Metodología}
Se emplea un aprendizaje basado en pequeños proyectos con trabajo cooperativo en grupos de 3-4 personas, con roles definidos y rotativos:
\begin{itemize}
    \item \textbf{Coordinador/a}: organiza el reparto de tareas y controla los plazos.
    \item \textbf{Documentalista}: busca y contrasta la información sobre la composición recomendada de los sueros de rehidratación.
    \item \textbf{Responsable de cálculo y de laboratorio}: realiza los cálculos de concentración y, con supervisión del profesorado, prepara la disolución de prueba (con sal y azúcar comestibles, sin componentes peligrosos).
    \item \textbf{Portavoz/diseñador/a}: elabora el folleto informativo final y lo presenta a la clase.
\end{itemize}
Se combina la búsqueda documental guiada, el cálculo matemático riguroso y una breve práctica de laboratorio supervisada, fomentando la autonomía, la responsabilidad individual dentro del grupo y la interacción verbal en la puesta en común final.

\subsection*{6. Recursos}
\begin{itemize}
    \item Balanza de precocina o de laboratorio, vasos de precipitados o recipientes graduados, agua, sal común y azúcar (materiales seguros y de bajo coste).
    \item Ordenadores o tablets del aula/dispositivo personal para la búsqueda de información.
    \item \textbf{Fuentes de información concretas a utilizar:}
    \begin{itemize}
        \item Página del \textbf{Ministerio de Sanidad} o de la \textbf{Organización Mundial de la Salud (OMS)} sobre soluciones de rehidratación oral.
        \item Prospecto o ficha técnica de un sobre de suero de rehidratación comercial (foto aportada por el alumnado o proporcionada por el profesorado).
        \item Buscador general (Google) con palabras clave sugeridas: \textit{``composición suero rehidratación oral OMS gramos por litro''}, \textit{``suero casero receta OMS sal azúcar''}.
    \end{itemize}
    \item \textbf{Criterio de fiabilidad de fuentes:} se priorizarán fuentes sanitarias oficiales (dominios .gob.es, who.int, sanidad.es) frente a blogs personales o foros no especializados; cada grupo debe justificar por escrito por qué la fuente elegida es fiable (organismo emisor, carácter oficial, fecha de publicación).
    \item \textbf{Formato mínimo de cita:} Organismo o autor. \textit{Título del documento o página}. Nombre del sitio web. URL completa. Fecha de consulta.
\end{itemize}

\subsection*{7. Tareas y actividades}
\begin{enumerate}
    \item \textbf{Fase 1 (grupo, 1 sesión):} Búsqueda documental guiada de la composición recomendada (en gramos por litro) de una solución de rehidratación oral casera, según fuentes sanitarias oficiales, citando correctamente la fuente utilizada.
    \item \textbf{Fase 2 (grupo, 1 sesión):} A partir de los datos encontrados, cálculo de la cantidad exacta de sal y de azúcar necesarias para preparar 500 mL y 1 L de la solución, expresando los resultados en porcentaje en masa y en concentración (g/L).
    \item \textbf{Fase 3 (grupo, 1 sesión, con supervisión del profesorado):} Preparación real de una pequeña cantidad de la disolución (500 mL) siguiendo los cálculos realizados, pesando los componentes con la balanza y verificando el procedimiento.
    \item \textbf{Fase 4 (grupo, 1 sesión):} Comparación de los datos calculados por el grupo con la composición de un sobre de suero comercial (a partir de su prospecto), analizando semejanzas y diferencias.
    \item \textbf{Fase 5 (grupo, 1 sesión):} Elaboración de un folleto informativo dirigido a las familias del centro, explicando qué es un suero de rehidratación, cuándo es recomendable usarlo y cómo prepararlo de forma segura en casa.
    \item \textbf{Fase 6 (gran grupo, 1 sesión):} Puesta en común y exposición oral breve de cada grupo (5 minutos), con turno de preguntas del resto de la clase.
\end{enumerate}

\textbf{Cuestiones que el alumnado deberá responder y entregar por escrito (individual o grupal) al finalizar:}
\begin{enumerate}
    \item ¿Qué cantidad de sal y de azúcar (en gramos) se necesita para preparar 1 litro de suero casero, según la fuente consultada? Indica la fuente y la fecha de consulta.
    \item Calcula el porcentaje en masa de sal y de azúcar en la disolución preparada de 500 mL, suponiendo una densidad de la disolución de 1,02 g/cm³.
    \item ¿Cuál es el soluto y cuál el disolvente en esta disolución? Justifica tu respuesta.
    \item Compara la concentración de sal (en g/L) de vuestra disolución casera con la que indica el prospecto de un sobre de suero comercial. ¿Son similares? ¿A qué puede deberse la diferencia, si la hay?
    \item Explica, a nivel de partículas, por qué al añadir la sal y el azúcar en agua y remover se obtiene una mezcla homogénea y no se pueden distinguir a simple vista los componentes.
    \item ¿Por qué es importante seguir las proporciones exactas al preparar un suero de rehidratación casero, y qué podría ocurrir si se pone demasiada sal?
    \item Reflexión grupal: ¿cómo os habéis repartido el trabajo? ¿Qué dificultad os ha resultado más complicada: la búsqueda de información fiable, los cálculos o la preparación experimental?
\end{enumerate}

\subsection*{8. Evaluación}
Se evaluará mediante rúbrica que valore: (a) la corrección de los cálculos de porcentaje en masa y concentración en masa (ligado a CE1.2 y CE3.2); (b) la fiabilidad y correcta citación de las fuentes de información sanitaria utilizadas (ligado a CE3.1 y CE4.2); (c) el rigor y la seguridad en la preparación experimental de la disolución (ligado a CE2.2); (d) la calidad y claridad del folleto informativo final, valorando su utilidad real para las familias (ligado a CE1.1); y (e) el funcionamiento cooperativo del grupo, a través de la reflexión final y la observación directa en el aula (ligado a CE5.1). Se combinará la evaluación del producto final con la observación continua del proceso.

\subsection*{9. Observaciones}
Es fundamental supervisar directamente la fase de preparación experimental para garantizar el uso seguro del material y evitar la ingesta de la disolución preparada (uso exclusivamente demostrativo/educativo, no para consumo real sin supervisión sanitaria). Se recomienda coordinar esta actividad con el departamento de Biología y Geología o con el equipo de enfermería escolar si el centro dispone de él, para reforzar el enfoque de educación para la salud. Puede ampliarse invitando a las familias a leer el folleto final en una reunión de tutoría o publicándolo en la web del centro.

%===================================================================
\newpage
\section{Unidad 6: Tabla Periódica y el Átomo}
\subsection*{1. Título}
``Detectives atómicos: investigando los elementos que nos rodean''

\subsection*{2. Justificación}
Esta situación de aprendizaje pretende que el alumnado consolide de forma significativa los conceptos de estructura atómica, isótopos, iones y organización de la tabla periódica trabajados en la unidad, mediante una investigación cooperativa sobre elementos químicos presentes en objetos y situaciones de la vida cotidiana. Se busca desarrollar la competencia científica (relacionar propiedades macroscópicas con la estructura atómica), la competencia digital (búsqueda y contraste de información) y la competencia de trabajo en equipo, mostrando que la tabla periódica no es un simple cartel de clase sino una herramienta que explica el mundo material que nos rodea.

\subsection*{3. Contextualización}
El alumnado convive a diario con objetos hechos de distintos elementos químicos sin ser consciente de ello: el aluminio del papel de cocina, el hierro de las herramientas, el cobre de los cables eléctricos, el litio de la batería del móvil, el yodo de la sal de mesa yodada, o el argón en las bombillas. Muchos de estos elementos, además, tienen una relevancia económica o industrial en Castilla-La Mancha (por ejemplo, la minería histórica del plomo y la plata en la región, o la industria del automóvil que usa aluminio y litio en baterías). Esta situación aprovecha ese contexto cercano para que el alumnado ``detective'' busque el elemento químico detrás de objetos y fenómenos cotidianos.

\subsection*{4. Fundamentación curricular}
\textbf{Competencias específicas relacionadas}:
\begin{itemize}
    \item \textbf{CE1} --- Comprender y explicar por qué las propiedades de los materiales cotidianos dependen de la estructura atómica y de la posición del elemento en la tabla periódica.
    \item \textbf{CE2} --- Emplear metodologías de indagación (uso de simuladores, búsqueda documental) para identificar propiedades atómicas.
    \item \textbf{CE3} --- Manejar con soltura la notación atómica, la nomenclatura de iones y la interpretación de datos de la tabla periódica.
    \item \textbf{CE4} --- Utilizar de forma crítica plataformas digitales (simuladores, tablas periódicas interactivas) para investigar propiedades de los elementos.
    \item \textbf{CE5} --- Trabajar de forma colaborativa en la elaboración de un producto final compartido con la clase.
    \item \textbf{CE6} --- Valorar la importancia histórica y social de la tabla periódica y de los elementos en la tecnología actual.
\end{itemize}

\textbf{Criterios de evaluación concretos}: CE1.1, CE1.2, CE2.1, CE3.1, CE3.2, CE4.2, CE5.1, CE6.1.

\textbf{Saberes básicos}:
\begin{itemize}
    \item \textbf{Bloque B (La materia)}: ``Estructura atómica: desarrollo histórico de los modelos atómicos, existencia, formación y propiedades de los isótopos y ordenación de los elementos en la tabla periódica.''
    \item \textbf{Bloque B (La materia)}: ``Nomenclatura: participación de un lenguaje científico común y universal formulando y nombrando sustancias simples, iones monoatómicos y compuestos binarios mediante las reglas de nomenclatura de la Iupac.''
    \item \textbf{Bloque A (Las destrezas científicas básicas)}: ``Diversos entornos y recursos de aprendizaje científico como el laboratorio o los entornos virtuales: materiales, sustancias y herramientas tecnológicas.''
    \item \textbf{Bloque A (Las destrezas científicas básicas)}: ``Valoración de la cultura científica y del papel de científicos y científicas en los principales hitos históricos y actuales de la física y la química en el avance y la mejora de la sociedad.''
\end{itemize}

\subsection*{5. Metodología}
Se emplea un enfoque de aprendizaje cooperativo basado en la investigación guiada (WebQuest simplificada), organizado en grupos heterogéneos de 3-4 personas, con roles definidos y rotativos:
\begin{itemize}
    \item \textbf{Coordinador/a}: organiza el reparto de tareas y controla los plazos del grupo.
    \item \textbf{Investigador/a histórico}: busca el descubrimiento y la historia del elemento asignado.
    \item \textbf{Investigador/a atómico}: usa el simulador PhET y calcula protones, neutrones, electrones, isótopos e iones relevantes del elemento.
    \item \textbf{Portavoz/diseñador/a}: elabora la ficha o cromo final del elemento y lo presenta a la clase.
\end{itemize}
Se combinan la manipulación de un simulador interactivo, la búsqueda documental contrastada y la puesta en común final, favoreciendo la autonomía, la interacción verbal y la construcción colectiva de una ``tabla periódica de aula''.

\subsection*{6. Recursos}
\begin{itemize}
    \item Ordenadores o tablets del aula/dispositivo personal.
    \item Simulador \textbf{PhET ``Build an Atom''} (Universidad de Colorado Boulder, acceso gratuito) para construir y comprobar la estructura atómica del elemento asignado.
    \item \textbf{Fuentes de información concretas a utilizar:}
    \begin{itemize}
        \item Tabla periódica interactiva oficial, como la de la \textbf{Royal Society of Chemistry (RSC)} o \textbf{ptable.com}.
        \item Buscador general (Google) con palabras clave sugeridas: \textit{``[nombre del elemento] descubrimiento historia''}, \textit{``[nombre del elemento] usos cotidianos propiedades''}.
        \item Enciclopedias o portales educativos de ciencia reconocidos (evitar páginas sin autoría).
    \end{itemize}
    \item \textbf{Criterio de fiabilidad de fuentes:} se priorizarán páginas de organismos científicos, universidades o editoriales educativas reconocidas, frente a blogs personales, foros o páginas sin autor ni fecha. Cada grupo debe justificar por qué considera fiable la fuente principal utilizada.
    \item \textbf{Formato mínimo de cita:} Autor/organismo (si consta). \textit{Título de la página o recurso}. Nombre del sitio web. URL completa. Fecha de consulta.
    \item Material para elaborar la ficha final: cartulina o plantilla digital de ``cromo de elemento'' (tipo carta coleccionable).
\end{itemize}

\subsection*{7. Tareas y actividades}
\begin{enumerate}
    \item \textbf{Fase 1 (grupo, 1 sesión):} Cada grupo recibe (o elige entre una lista propuesta) un elemento químico presente en un objeto cotidiano: aluminio, hierro, cobre, litio, yodo, argón, oro, silicio, etc.
    \item \textbf{Fase 2 (grupo, 1 sesión):} Con el simulador PhET ``Build an Atom'', el grupo construye el átomo neutro del elemento asignado, comprobando el número de protones, neutrones y electrones, y explora qué ocurre al formar un ion (catión o anión) típico de ese elemento.
    \item \textbf{Fase 3 (grupo, 1 sesión):} Cálculo en papel de la configuración electrónica del elemento, su valencia (electrones de valencia) y su localización (grupo y periodo) en la tabla periódica, contrastando el resultado con una tabla periódica real.
    \item \textbf{Fase 4 (grupo, 1 sesión):} Búsqueda documental guiada sobre la historia del descubrimiento del elemento y sus usos cotidianos o industriales (incluyendo, si aplica, su relevancia en Castilla-La Mancha), citando correctamente las fuentes utilizadas.
    \item \textbf{Fase 5 (grupo, 1 sesión):} Elaboración de una ficha o ``cromo coleccionable'' del elemento con: símbolo, Z, A, configuración electrónica, grupo/periodo, un dato histórico y un uso cotidiano, ilustrado con un dibujo o imagen.
    \item \textbf{Fase 6 (gran grupo, 1 sesión):} Presentación oral breve de cada grupo (4-5 minutos) y montaje conjunto de una ``tabla periódica de aula'' con todas las fichas elaboradas, colocadas en su grupo y periodo correspondiente.
\end{enumerate}

\textbf{Cuestiones que el alumnado deberá responder y entregar por escrito (individual o grupal) al finalizar:}
\begin{enumerate}
    \item ¿Cuántos protones, neutrones y electrones tiene un átomo neutro de vuestro elemento? Indicad Z y A.
    \item Escribid la configuración electrónica completa de vuestro elemento y señalad su capa de valencia y el número de electrones de valencia.
    \item ¿En qué grupo y periodo de la tabla periódica se encuentra vuestro elemento? ¿A qué familia pertenece (metal alcalino, halógeno, gas noble, metal de transición, etc.)?
    \item Si vuestro elemento forma habitualmente un ion, indicad su carga y el número de electrones que tiene ese ion. Explicad por qué tiende a formar ese ion en concreto, relacionándolo con sus electrones de valencia.
    \item ¿Tiene vuestro elemento isótopos conocidos? Si es así, indicad al menos uno y explicad en qué se diferencia del isótopo más común.
    \item ¿Quién descubrió vuestro elemento y en qué año, según la fuente consultada? ¿Qué uso cotidiano o industrial tiene hoy en día?
    \item Reflexión grupal: ¿cómo os habéis repartido el trabajo de investigación? ¿Qué os ha resultado más difícil, el cálculo atómico o la búsqueda de información histórica fiable?
\end{enumerate}

\subsection*{8. Evaluación}
Se evaluará mediante rúbrica que contemple: (a) la corrección de los cálculos de protones/neutrones/electrones y de la configuración electrónica (ligado a CE1.2 y CE3.2); (b) el uso correcto del simulador PhET como herramienta de comprobación experimental (ligado a CE2.1); (c) la correcta localización del elemento en la tabla periódica y la interpretación de su grupo/familia (ligado a CE3.1 y CE1.1); (d) la fiabilidad y correcta citación de las fuentes históricas utilizadas (ligado a CE4.2 y CE6.1); (e) la calidad y creatividad de la ficha/cromo final; y (f) el funcionamiento cooperativo del grupo, valorado mediante la reflexión final escrita y la observación directa del profesorado durante las sesiones (ligado a CE5.1). Se combinará la evaluación del producto final con la observación continua del proceso de trabajo.

\subsection*{9. Observaciones}
Es recomendable repartir los elementos entre los grupos evitando duplicados y asegurando una variedad representativa de familias (metales, no metales, gases nobles) para que la ``tabla periódica de aula'' final quede lo más completa posible. Conviene reservar un aula con ordenadores o permitir el uso de dispositivos personales para garantizar el acceso al simulador PhET. La tabla periódica de aula resultante puede quedar expuesta en el pasillo del centro como recurso visual para otros grupos, dando proyección al trabajo realizado.

%===================================================================
\newpage
\section{Nomenclatura Inorgánica}
\subsection*{1. Título}
Lo que dice la etiqueta: químicos seguros en casa y en el campo de Castilla-La Mancha

\subsection*{2. Justificación}
El alumnado convive a diario con productos etiquetados con nombres y símbolos químicos que no comprende: productos de limpieza, fertilizantes y fitosanitarios usados en la agricultura de la región (muy presente en las familias de zonas rurales de Castilla-La Mancha), o medicamentos de uso común. Esta situación de aprendizaje pretende que el alumnado aplique las reglas de nomenclatura y formulación inorgánica aprendidas en la unidad para \textbf{interpretar de forma crítica y segura} la información química de su entorno cercano, desarrollando a la vez la competencia digital (búsqueda de información fiable) y el trabajo en equipo.

\subsection*{3. Contextualización}
En muchas localidades de Castilla-La Mancha conviven el uso doméstico de productos de limpieza (lejía, salfumán, agua oxigenada) con el uso agrícola de fertilizantes (sulfato de amonio, superfosfatos) y correctores de suelo como la cal agrícola (\ch{Ca(OH)2}), habitual en el manejo de los suelos calizos de la región. Cada año se producen accidentes domésticos evitables por mezclar productos incompatibles (p. ej. lejía y salfumán). El reto se plantea así: \textit{``El ayuntamiento de tu localidad quiere publicar un folleto informativo para las familias sobre el uso seguro de los productos químicos más comunes en casa y en el huerto. Tu equipo ha sido contratado como consultores químicos para elaborarlo''}.

\subsection*{4. Fundamentación curricular}
\textbf{Competencias específicas}:
\begin{itemize}
    \item \textbf{CE1} --- Comprensión y resolución de fenómenos fisicoquímicos cotidianos.
    \item \textbf{CE3} --- Manejo del lenguaje científico (IUPAC, unidades, seguridad en el laboratorio).
    \item \textbf{CE4} --- Uso crítico y seguro de plataformas digitales y recursos variados.
    \item \textbf{CE5} --- Trabajo colaborativo.
\end{itemize}

\textbf{Criterios de evaluación trabajados}: CE1.1, CE1.3, CE3.2, CE3.3, CE4.1, CE4.2, CE5.1, CE5.2.

\textbf{Saberes básicos} (bloque \textbf{B. La materia}):
\begin{itemize}
    \item \textit{Principales compuestos químicos: su formación y sus propiedades físicas y químicas, valoración de sus aplicaciones. Masa atómica y masa molecular.}
    \item \textit{Nomenclatura: participación de un lenguaje científico común y universal formulando y nombrando sustancias simples, iones monoatómicos y compuestos binarios mediante las reglas de nomenclatura de la Iupac.}
\end{itemize}
(bloque \textbf{A. Las destrezas científicas básicas}):
\begin{itemize}
    \item \textit{Normas de uso de cada espacio, asegurando y protegiendo así la salud propia y comunitaria, la seguridad en las redes y el respeto hacia el medio ambiente.}
    \item \textit{Estrategias de interpretación y producción de información científica utilizando diferentes formatos y diferentes medios.}
\end{itemize}

\subsection*{5. Metodología}
Se emplea \textbf{aprendizaje cooperativo basado en proyectos} (ABP). El alumnado se organiza en \textbf{grupos de 3-4 personas}, con roles rotativos y definidos:
\begin{itemize}
    \item \textbf{Coordinador/a}: organiza el reparto de tareas y los plazos.
    \item \textbf{Investigador/a químico/a}: busca y verifica nombres, fórmulas y propiedades de los compuestos.
    \item \textbf{Responsable de seguridad}: investiga pictogramas de riesgo y normas de manipulación.
    \item \textbf{Editor/a-portavoz}: redacta y maqueta el folleto final y lo presenta al resto de la clase.
\end{itemize}
El profesorado actúa como guía, resolviendo dudas y validando las fuentes elegidas, favoreciendo la autonomía del alumnado y la interacción verbal en el grupo.

\subsection*{6. Recursos}
\begin{itemize}
    \item Etiquetas reales (o fotografías tomadas por el alumnado, con permiso familiar) de productos de casa: lejía, salfumán, agua oxigenada, quitamanchas, fertilizante de jardín.
    \item Base de datos química \textbf{PubChem} (\url{https://pubchem.ncbi.nlm.nih.gov}) para buscar el nombre IUPAC y la fórmula de un compuesto a partir de su nombre comercial (palabras clave sugeridas: nombre del principio activo en inglés, p.ej. \textit{``sodium hypochlorite''}, \textit{``hydrogen peroxide''}).
    \item Fichas de datos de seguridad (FDS) oficiales, buscables mediante el buscador del centro con las palabras clave \textit{``ficha de datos de seguridad'' + nombre del producto}.
    \item Pictogramas de peligrosidad GHS (sistema globalmente armonizado), disponibles en la web de la Agencia Europea de Sustancias y Mezclas Químicas (ECHA, \url{https://echa.europa.eu}).
    \item Apuntes de nomenclatura de la unidad y tabla de estados de oxidación del libro de texto.
    \item Procesador de texto o herramienta de diseño de folletos (Canva, LibreOffice) para el producto final.
\end{itemize}
\textbf{Criterio de fiabilidad de fuentes}: solo se aceptan fuentes institucionales u oficiales (bases de datos científicas, agencias gubernamentales o europeas, fabricantes con ficha de seguridad oficial); no se aceptan foros o blogs sin autoría verificable. \textbf{Formato de cita mínimo}: nombre de la fuente, URL y fecha de consulta, anotados en un apartado de ``Fuentes consultadas'' al final del folleto (ej.: \textit{PubChem, https://pubchem.ncbi.nlm.nih.gov/compound/..., consultado el 11/03/2026}).

\subsection*{7. Tareas y actividades}
\begin{enumerate}
    \item Cada grupo recibe (o fotografía) 4 productos químicos reales de uso doméstico o agrícola.
    \item Para cada producto: identificar el/los compuesto(s) químico(s) activo(s), obtener su fórmula y su nombre según la nomenclatura de estados de oxidación, usando PubChem y contrastando con los apuntes de la unidad.
    \item Buscar el pictograma de peligro correspondiente en ECHA y redactar una norma de manipulación segura para cada producto.
    \item Calcular la masa molecular de al menos dos de los compuestos identificados.
    \item Elaborar el folleto final con la información recogida, incluyendo el apartado de fuentes consultadas.
    \item Presentación oral de 5 minutos por grupo al resto de la clase.
\end{enumerate}
\textbf{Cuestiones que el alumnado deberá responder y entregar} (trabajo escrito grupal, junto al folleto):
\begin{enumerate}
    \item ¿Cuál es el nombre químico (multiplicadores y estado de oxidación) y la fórmula de cada uno de los 4 productos investigados?
    \item ¿Qué tipo de compuesto inorgánico es cada uno (óxido, hidróxido, sal binaria, hidruro, ácido...)?
    \item ¿Cuál es la masa molecular de dos de los compuestos elegidos y cómo se ha calculado?
    \item ¿Qué pictograma(s) de peligro lleva cada producto y qué significan?
    \item ¿Por qué es peligroso mezclar lejía (hipoclorito de sodio) y salfumán (ácido clorhídrico)? Justifica la respuesta usando lo aprendido sobre reacciones y compuestos.
    \item ¿Qué fuentes habéis utilizado para obtener la información y por qué las consideráis fiables?
    \item Reflexión personal: ¿qué producto de los estudiados desconocías que tenía ese nombre químico y qué te ha sorprendido más?
\end{enumerate}

\subsection*{8. Evaluación}
Se evaluará mediante rúbrica grupal e individual, atendiendo a los criterios CE1.1, CE1.3, CE3.2, CE3.3, CE4.1, CE4.2, CE5.1 y CE5.2:
\begin{itemize}
    \item \textbf{Corrección científica} (nombres, fórmulas, masas moleculares): CE3.2, CE1.1.
    \item \textbf{Tratamiento de la seguridad y aplicación a la vida real}: CE1.3, CE3.3.
    \item \textbf{Uso y cita adecuada de fuentes digitales fiables}: CE4.1, CE4.2.
    \item \textbf{Funcionamiento del equipo y desempeño del rol asignado}: CE5.1, CE5.2 (mediante coevaluación entre miembros del grupo y observación del profesorado).
    \item Las respuestas escritas entregadas se corrigen de forma individual dentro de la nota grupal, para asegurar la implicación de cada estudiante.
\end{itemize}

\subsection*{9. Observaciones}
Se recomienda avisar a las familias con antelación si se solicita traer fotografías de productos de casa, por motivos de privacidad y seguridad (nunca se debe manipular lejía y salfumán juntos ni traerlos físicamente al aula). Es aconsejable programar esta situación de aprendizaje al finalizar la unidad de nomenclatura, como cierre práctico y transferencia de lo aprendido, dedicándole 3 sesiones de 50 minutos (búsqueda e investigación, elaboración del folleto, exposición final).

%===================================================================
\newpage
\section{Unidad 7: Elementos y Compuestos}
\subsection*{1. Título}
¿Cuántas moléculas de agua se pierden? Investigando el ahorro de agua en Castilla-La Mancha con el concepto de mol

\subsection*{2. Justificación}
El concepto de mol resulta muy abstracto para el alumnado de 3º de ESO. Esta situación de aprendizaje pretende dar significado real a ese concepto conectándolo con un problema muy presente en Castilla-La Mancha: la escasez de agua y la necesidad de un uso responsable del recurso hídrico. El alumnado aplicará los cálculos de masa molar, mol y número de partículas a datos reales de consumo de agua, desarrollando a la vez conciencia ambiental, competencia digital y trabajo en equipo.

\subsection*{3. Contextualización}
Castilla-La Mancha es una de las regiones españolas más afectadas por la escasez de agua, con episodios de sequía recurrentes que afectan a los embalses y acuíferos de la región (como el acuífero 23 de la Mancha) y condicionan la agricultura y el consumo doméstico. El reto se plantea así: \textit{``El ayuntamiento de tu localidad ha pedido a los centros educativos que preparen una campaña de concienciación sobre el ahorro de agua, mostrando de una forma sorprendente y rigurosa cuánta 'cantidad de sustancia' (en moles y en número de moléculas) se gasta o se ahorra con nuestros hábitos diarios''}.

\subsection*{4. Fundamentación curricular}
\textbf{Competencias específicas}:
\begin{itemize}
    \item \textbf{CE1} --- Comprensión y resolución de fenómenos fisicoquímicos cotidianos.
    \item \textbf{CE3} --- Manejo del lenguaje científico (unidades, datos, mol).
    \item \textbf{CE4} --- Uso crítico y seguro de plataformas digitales y recursos variados.
    \item \textbf{CE5} --- Trabajo colaborativo.
    \item \textbf{CE6} --- La ciencia y su repercusión social y ambiental.
\end{itemize}

\textbf{Criterios de evaluación trabajados}: CE1.1, CE1.2, CE3.1, CE4.1, CE4.2, CE5.1, CE5.2, CE6.2.

\textbf{Saberes básicos}:
\begin{itemize}
    \item (bloque \textbf{B. La materia}) \textit{Principales compuestos químicos: su formación y sus propiedades físicas y químicas, valoración de sus aplicaciones. Masa atómica y masa molecular.}
    \item (bloque \textbf{A. Las destrezas científicas básicas}) \textit{El lenguaje científico: unidades del Sistema Internacional y sus símbolos. Herramientas matemáticas básicas en diferentes escenarios científicos y de aprendizaje.}
    \item (bloque \textbf{C. La energía}) \textit{Elaboración fundamentada de hipótesis sobre el medio ambiente y la sostenibilidad...} (se adapta aquí al recurso hídrico como recurso natural gestionado de forma sostenible).
\end{itemize}

\subsection*{5. Metodología}
Aprendizaje cooperativo basado en un reto (ABP), en \textbf{grupos de 3-4 estudiantes} con roles definidos y rotativos:
\begin{itemize}
    \item \textbf{Coordinador/a}: organiza el trabajo y los plazos del grupo.
    \item \textbf{Analista de datos}: busca y contrasta los datos reales de consumo de agua.
    \item \textbf{Calculista}: realiza y verifica los cálculos de masa, moles y partículas.
    \item \textbf{Comunicador/a}: diseña el material de la campaña y presenta los resultados a la clase.
\end{itemize}
El profesorado guía el proceso de indagación, resuelve dudas técnicas sobre las conversiones y valida las fuentes de datos utilizadas, fomentando la autonomía y la interacción verbal dentro del grupo.

\subsection*{6. Recursos}
\begin{itemize}
    \item Simulador \textbf{PhET ``Molarity''} y \textbf{PhET ``States of Matter''} (\url{https://phet.colorado.edu}) para visualizar partículas y cantidad de sustancia.
    \item Datos reales de consumo de agua: web de la \textbf{Confederación Hidrográfica del Guadiana/Júcar} o del Instituto Nacional de Estadística (INE), buscados con palabras clave como \textit{``consumo doméstico de agua por habitante Castilla-La Mancha''}, \textit{``embalses Castilla-La Mancha datos''}.
    \item Tabla periódica interactiva (p. ej. \url{https://ptable.com}) para consultar masas atómicas.
    \item Calculadora científica y hoja de cálculo (LibreOffice Calc o similar) para organizar los cálculos.
    \item Apuntes de la unidad (mol, masa molar, número de Avogadro).
\end{itemize}
\textbf{Criterio de fiabilidad de fuentes}: se exige el uso de fuentes oficiales (organismos públicos, confederaciones hidrográficas, INE) o científicas (PubChem, PhET, ptable.com); no se aceptan blogs personales o webs sin autoría identificable. \textbf{Formato de cita mínimo}: nombre del organismo/fuente, URL y fecha de consulta (ej.: \textit{INE, https://www.ine.es/..., consultado el 10/03/2026}).

\subsection*{7. Tareas y actividades}
\begin{enumerate}
    \item Buscar el dato real de consumo medio de agua por persona y día en la región (litros/persona/día).
    \item Calcular, a partir de ese dato, cuántos moles de agua (\ch{H2O}) y cuántas moléculas de agua consume una persona en un día, en una semana y en un año.
    \item Calcular cuántas moléculas de agua se ahorrarían con un gesto cotidiano concreto (p. ej. cerrar el grifo al lavarse los dientes, ducha de 5 min en vez de 10 min), buscando o estimando de forma razonada el volumen de agua ahorrado.
    \item Diseñar un cartel o infografía de campaña con los resultados, pensada para el público general (sin lenguaje demasiado técnico, pero con los datos correctos).
    \item Presentación oral de 5 minutos por grupo mostrando los cálculos y la propuesta de campaña.
\end{enumerate}
\textbf{Cuestiones que el alumnado deberá responder y entregar} (trabajo escrito grupal):
\begin{enumerate}
    \item ¿Cuál es la fuente de datos utilizada para el consumo de agua y por qué la consideráis fiable?
    \item ¿Cuántos moles y cuántas moléculas de agua consume una persona al día, según vuestros cálculos? Mostrad el procedimiento completo.
    \item ¿Cuántas moléculas de agua se ahorrarían al año con el gesto cotidiano que habéis elegido?
    \item Explicad, con vuestras palabras, qué relación existe entre litros, gramos, moles y moléculas, y por qué es necesario pasar por los moles para llegar de litros a moléculas.
    \item ¿Por qué creéis que el ahorro de agua es especialmente importante en Castilla-La Mancha? Aportad al menos un dato o argumento encontrado en vuestra investigación.
    \item Reflexión personal: ¿cambia algo tu percepción del ``gasto de agua'' al expresarlo en número de moléculas en lugar de en litros? Justifica tu respuesta.
\end{enumerate}

\subsection*{8. Evaluación}
Rúbrica que valora, en relación con los criterios trabajados:
\begin{itemize}
    \item \textbf{Corrección de los cálculos} (masa, moles, partículas): CE1.2, CE3.1.
    \item \textbf{Explicación conceptual de los fenómenos y magnitudes}: CE1.1.
    \item \textbf{Uso y cita adecuada de fuentes fiables}: CE4.1, CE4.2.
    \item \textbf{Funcionamiento del equipo y desempeño de los roles}: CE5.1, CE5.2 (coevaluación grupal + observación docente).
    \item \textbf{Argumentación sobre la relevancia social/ambiental del reto}: CE6.2.
    \item Las respuestas escritas se corrigen de forma individualizada dentro de la nota grupal para garantizar la implicación de cada estudiante.
\end{itemize}

\subsection*{9. Observaciones}
Se recomienda coordinar esta situación de aprendizaje con el departamento de Biología y Geología o con el plan de sostenibilidad del centro, si existe, para reforzar la transversalidad del reto. Duración estimada: 3 sesiones de 50 minutos (investigación de datos y cálculo, elaboración del material de campaña, exposición final). Es recomendable exponer los carteles resultantes en un espacio visible del centro como acción real de sensibilización.

%===================================================================
\newpage
\section{Unidad 8: Reacciones Químicas}
\subsection*{1. Título}
De la cantera a la casa: la química de la cal en Castilla-La Mancha

\subsection*{2. Justificación}
Las reacciones químicas de descomposición y combinación se estudian a menudo de forma abstracta, alejadas de la experiencia del alumnado. Esta situación de aprendizaje conecta la reacción de calcinación de la piedra caliza ($\ch{CaCO3} \rightarrow \ch{CaO} + \ch{CO2}$) con una actividad económica real e histórica de numerosas localidades de Castilla-La Mancha (hornos de cal, construcción tradicional con cal y yeso), permitiendo al alumnado aplicar el ajuste de ecuaciones y los cálculos estequiométricos a un proceso productivo real, desarrollando además conciencia sobre el impacto ambiental de este tipo de procesos industriales.

\subsection*{3. Contextualización}
En Castilla-La Mancha existe una larga tradición de fabricación de cal y yeso (hornos de cal en numerosas comarcas de Cuenca, Toledo o Guadalajara), un material todavía usado hoy en restauración de edificios tradicionales y en agricultura (encalado de suelos). El reto se plantea así: \textit{``Una empresa local de materiales de construcción quiere modernizar su producción de cal y os pide, como jóvenes asesores científicos, un informe técnico que explique la reacción química implicada, calcule cuánta cal y cuánto \ch{CO2} se produce a partir de una cantidad de piedra caliza, y valore el impacto ambiental del proceso''}.

\subsection*{4. Fundamentación curricular}
\textbf{Competencias específicas}:
\begin{itemize}
    \item \textbf{CE1} --- Comprensión y resolución de fenómenos fisicoquímicos cotidianos.
    \item \textbf{CE3} --- Manejo del lenguaje científico y de datos en distintos formatos.
    \item \textbf{CE4} --- Uso crítico y seguro de plataformas digitales y recursos variados.
    \item \textbf{CE5} --- Trabajo colaborativo.
    \item \textbf{CE6} --- La ciencia como construcción colectiva y su repercusión social y ambiental.
\end{itemize}

\textbf{Criterios de evaluación trabajados}: CE1.1, CE1.2, CE3.1, CE4.1, CE4.2, CE5.1, CE5.2, CE6.2.

\textbf{Saberes básicos}:
\begin{itemize}
    \item (bloque \textbf{E. El cambio}) \textit{Interpretación macroscópica y microscópica de las reacciones químicas: explicación de las relaciones de la química con el medio ambiente, la tecnología y la sociedad.}
    \item (bloque \textbf{E. El cambio}) \textit{Ley de conservación de la masa y de la ley de las proporciones definidas: aplicación de estas leyes como evidencias experimentales que permiten validar el modelo atómico-molecular de la materia.}
    \item (bloque \textbf{E. El cambio}) \textit{Factores que afectan a las reacciones químicas: predicción cualitativa de la evolución de las reacciones, entendiendo su importancia en la resolución de problemas actuales por parte de la ciencia.}
    \item (bloque \textbf{A. Las destrezas científicas básicas}) \textit{El lenguaje científico: unidades del Sistema Internacional y sus símbolos. Herramientas matemáticas básicas en diferentes escenarios científicos y de aprendizaje.}
\end{itemize}

\subsection*{5. Metodología}
Aprendizaje cooperativo basado en un reto real (ABP), en \textbf{grupos de 3-4 estudiantes}, con roles definidos y rotativos:
\begin{itemize}
    \item \textbf{Coordinador/a}: organiza el trabajo, controla los plazos.
    \item \textbf{Investigador/a de proceso industrial}: busca información sobre el proceso real de calcinación y los hornos de cal en la región.
    \item \textbf{Calculista}: realiza y verifica el ajuste de la ecuación y los cálculos estequiométricos.
    \item \textbf{Responsable de impacto ambiental y comunicación}: investiga las emisiones de \ch{CO2} asociadas y redacta el informe final.
\end{itemize}
El profesorado actúa como guía técnico, resolviendo dudas sobre el ajuste de ecuaciones y los cálculos, y validando la fiabilidad de las fuentes utilizadas, favoreciendo la autonomía y la interacción verbal dentro del grupo.

\subsection*{6. Recursos}
\begin{itemize}
    \item Simulador \textbf{PhET ``Balancing Chemical Equations''} (\url{https://phet.colorado.edu}) para practicar y verificar el ajuste de la reacción de calcinación y otras relacionadas.
    \item Buscador de información sobre patrimonio industrial: palabras clave sugeridas \textit{``hornos de cal Castilla-La Mancha''}, \textit{``industria de la cal patrimonio''}, para localizar información sobre hornos de cal históricos de la región (páginas de ayuntamientos, diputaciones o de patrimonio industrial).
    \item Datos sobre emisiones de \ch{CO2} de la industria de la cal, buscables con palabras clave como \textit{``emisiones CO2 industria de la cal''} en fuentes institucionales (Ministerio para la Transición Ecológica, agencias ambientales europeas).
    \item Tabla periódica interactiva (\url{https://ptable.com}) para las masas atómicas.
    \item Apuntes de la unidad (ajuste de ecuaciones, ley de conservación de la masa, cálculos estequiométricos).
\end{itemize}
\textbf{Criterio de fiabilidad de fuentes}: solo se aceptan fuentes institucionales, académicas o de patrimonio oficial (ayuntamientos, diputaciones, ministerios, agencias ambientales, simuladores educativos reconocidos); no se aceptan foros o blogs sin autoría verificable. \textbf{Formato de cita mínimo}: nombre de la fuente/organismo, URL y fecha de consulta, recogidos en un apartado final de ``Fuentes consultadas'' del informe (ej.: \textit{Ministerio para la Transición Ecológica, https://www.miteco.gob.es/..., consultado el 12/03/2026}).

\subsection*{7. Tareas y actividades}
\begin{enumerate}
    \item Investigar brevemente el proceso tradicional/industrial de calcinación de la piedra caliza y su presencia histórica en Castilla-La Mancha.
    \item Escribir y ajustar la ecuación química de la calcinación: $\ch{CaCO3} \rightarrow \ch{CaO} + \ch{CO2}$, verificando el ajuste con el simulador PhET.
    \item Calcular, a partir de una cantidad de piedra caliza asignada al grupo (p. ej. \qty{500}{\kg}), la masa de cal (\ch{CaO}) obtenida y la masa y el volumen (a STP) de \ch{CO2} liberado.
    \item Investigar el dato aproximado de emisiones de \ch{CO2} asociadas a la industria de la cal y compararlo con el resultado de vuestro cálculo.
    \item Redactar el informe técnico final con los cálculos, la investigación histórica y una valoración del impacto ambiental.
    \item Presentación oral de 5 minutos por grupo a la clase, actuando como ``asesores científicos'' ante la empresa ficticia del reto.
\end{enumerate}
\textbf{Cuestiones que el alumnado deberá responder y entregar} (trabajo escrito grupal, dentro del informe técnico):
\begin{enumerate}
    \item ¿Cuál es la ecuación química ajustada de la calcinación de la piedra caliza? Mostrad el proceso de ajuste paso a paso.
    \item Si se calcinan \qty{500}{\kg} de \ch{CaCO3}, ¿qué masa de \ch{CaO} y qué masa de \ch{CO2} se obtienen? Mostrad todos los cálculos ($\ch{Ca}=40$, $\ch{C}=12$, $\ch{O}=16$).
    \item ¿Qué volumen ocupa ese \ch{CO2} en condiciones estándar (STP)?
    \item ¿Se cumple la ley de conservación de la masa en vuestros cálculos? Comprobadlo numéricamente.
    \item ¿Qué información habéis encontrado sobre los hornos de cal históricos en Castilla-La Mancha? Citad la fuente utilizada.
    \item ¿Qué impacto ambiental tiene la liberación de \ch{CO2} en este proceso y qué medida propondríais para reducirlo o compensarlo?
    \item Reflexión personal: ¿qué has aprendido sobre la relación entre una reacción química ``de libro'' y un proceso industrial real de tu región?
\end{enumerate}

\subsection*{8. Evaluación}
Rúbrica que valora, en relación con los criterios trabajados:
\begin{itemize}
    \item \textbf{Corrección del ajuste de la ecuación y de los cálculos estequiométricos}: CE1.2, CE3.1.
    \item \textbf{Explicación conceptual de la reacción y de la ley de conservación de la masa}: CE1.1.
    \item \textbf{Uso y cita adecuada de fuentes fiables}: CE4.1, CE4.2.
    \item \textbf{Funcionamiento del equipo y desempeño de los roles asignados}: CE5.1, CE5.2 (coevaluación grupal + observación docente).
    \item \textbf{Argumentación sobre el impacto ambiental y la relevancia social del proceso}: CE6.2.
    \item Las respuestas escritas del informe se corrigen de forma individualizada dentro de la nota grupal para garantizar la implicación de cada estudiante.
\end{itemize}

\subsection*{9. Observaciones}
Se recomienda, si es posible, contactar con una empresa local de materiales de construcción o con el área de patrimonio del ayuntamiento para una charla breve o visita virtual, reforzando la conexión con la realidad del reto. Duración estimada: 3 sesiones de 50 minutos (investigación y ajuste de la ecuación, cálculos estequiométricos y redacción del informe, exposición final). Se recomienda situar esta situación de aprendizaje tras haber trabajado el ajuste de ecuaciones y antes de cerrar la unidad, como aplicación integradora de todo lo aprendido.

\end{document}
